% concatenated from UnivNorms.tex
\documentclass[12pt,english,a4paper]{smfart}

\usepackage[T1]{fontenc}
\usepackage{lmodern}
\usepackage{smfthm}
\usepackage[headings]{fullpage}

\newcommand{\Qp}{\mathbf{Q}_p}
\newcommand{\Zp}{\mathbf{Z}_p}
\newcommand{\ZZ}{\mathbf{Z}}
\renewcommand{\phi}{\varphi}
\renewcommand{\projlim}{\varprojlim}
\newcommand{\Gal}{\mathrm{Gal}}
\newcommand{\Iw}{\mathrm{Iw}}
\newcommand{\Fil}{\mathrm{Fil}}
\newcommand{\Nm}{\mathrm{N}}
\newcommand{\alg}{\mathrm{alg}}
\newcommand{\cyc}{\mathrm{cyc}}
\newcommand{\dfont}{\mathrm{D}}
\newcommand{\pr}{\operatorname{pr}}

\author{Laurent Berger}
\address{UMPA ENS de Lyon \\
UMR 5669 du CNRS}
\email{laurent.berger@ens-lyon.fr}
\urladdr{perso.ens-lyon.fr/laurent.berger/}

\title{Universal norms and de Rham representations}

\begin{document}

\begin{abstract}
The purpose of this note is to show that the computation of the universal norms in Iwasawa cohomology for the $p$-adic cyclotomic extension implies the same result for any $p$-adic extension that contains the cyclotomic extension. In addition, we prove a twisted cyclotomic variant, which applies to Lubin--Tate extensions.
\end{abstract}

\date{\today}

\thanks{The author's research was partially supported by the ANR grant ANR-25-CE40-4664}

\maketitle

\setlength{\baselineskip}{17pt}

\section{Main result}

Let $K$ be a finite extension of $\Qp$ (or, more generally, a finite extension of $W(k)[1/p]$, where $k$ is a perfect field of characteristic $p$). Let $G_K = \Gal(K^{\alg}/K)$ denote the absolute Galois group of $K$ and let $V$ be a $p$-adic representation of $G_K$. Let $K_\infty \subset K^{\alg}$ be an infinite extension of $K$. We consider the Iwasawa cohomology group $H^1_{\Iw}(K_\infty/K,V) = \Qp \otimes_{\Zp} \projlim_L H^1(L,T)$ where $T$ is a lattice of $V$, $L$ runs through the set $E(K_\infty/K)$ of finite subextensions of $K_\infty/K$, and the transition maps are the corestriction maps. If $V$ is a de Rham representation, we consider the submodule $H^1_{\Iw,g}(K_\infty/K,V)$ of $H^1_{\Iw}(K_\infty/K,V)$ consisting of those classes whose image in $H^1(L,V)$ is in the subspace $H^1_g(L,V)$ of $H^1(L,V)$ parametrizing de Rham extensions, for every $L \in E(K_\infty/K)$. 

Let $\Fil^1 V$ be the largest subrepresentation of $V$ whose Hodge--Tate weights are all $\geq 1$. The space $H^1_{\Iw}(K_\infty/K,\Fil^1 V)$ is then contained in $H^1_{\Iw,g}(K_\infty/K,V)$. When $K_\infty = K_{\cyc} := K( \mu_{p^\infty} )$ is the cyclotomic extension of $K$, the main result of \cite{Ber05} is that this inclusion is almost an equality. More precisely, Theorem A of ibid is the following 
(where $H_K = \Gal(K^{\alg}/K_{\cyc})$).

\begin{theo}
\label{theoA}
Let $V$ be a de Rham representation of $G_K$ and let $K_{\cyc}/K$ be the cyclotomic extension.
\begin{enumerate}
\item If $V$ has no subquotient fixed by $H_K$, then $H^1_{\Iw,g}(K_{\cyc}/K,V) = H^1_{\Iw}(K_{\cyc}/K,\Fil^1 V) $
\item In general, $H^1_{\Iw}(K_{\cyc}/K,\Fil^1 V) \subset H^1_{\Iw,g}(K_{\cyc}/K,V)$ and the quotient is a finite-dimensional $\Qp$-vector space.
\end{enumerate}
\end{theo}

This result proved a conjecture of Nekov\'a\v{r} (conjecture A of \cite{PR00}), building on earlier work of Perrin-Riou \cite{PR00} who used her big exponential map. The proof of Theorem \ref{theoA} uses the theory of $(\phi,\Gamma)$-modules. For other extensions $K_\infty/K$, in the absence of a good theory of $(\phi,\Gamma)$-modules, proving an analogue of this result seems difficult. There has however been recent encouraging progress by Ponsinet \cite{Pon24,Pon25}, using the Fargues--Fontaine curve instead of $(\phi,\Gamma)$-modules. 

The purpose of this note is to show that Theorem \ref{theoA}, together with the methods used to prove it, implies the following result.

\begin{theo}
\label{theoB}
If $V$ is a de Rham representation of $G_K$ and if $K_\infty$ \emph{contains} the cyclotomic extension $K_{\cyc}$ of $K$, then $H^1_{\Iw,g}(K_\infty/K,V)  / H^1_{\Iw}(K_\infty/K,\Fil^1 V)$ is a finite-dimensional $\Qp$-vector space.
\end{theo}

The idea behind this theorem is quite simple. Recall that $E(K_\infty/K)$ denotes the set of all finite subextensions $L$ of $K_\infty / K$. First suppose, for simplicity, that $V$ is such that for all $L \in E(K_\infty/K)$, we have $H^1_{\Iw,g}(L_{\cyc}/L,V) = H^1_{\Iw}(L_{\cyc}/L,\Fil^1 V)$ (namely that for all $L$, we are in case (1) of Theorem \ref{theoA}). Take $c \in H^1_{\Iw,g}(K_\infty/K,V)$ and for $n \geq 0$, let $L_n = L(\mu_{p^n})$ so that $L_{\cyc} = \cup_{n \geq 0} L_n$. For every such $L$ and every $n \geq 0$, we have $\pr_{L_n}(c) \in H^1_g(L_n,V)$. Thus, the compatible collection of classes $\{ \pr_{L_n}(c) \}_{n \geq 0}$ defines an element of $H^1_{\Iw,g}(L_{\cyc}/L,V)$. Since we are assuming that $H^1_{\Iw,g}(L_{\cyc}/L,V) = H^1_{\Iw}(L_{\cyc}/L,\Fil^1 V)$, we have $\pr_L(c) \in H^1(L, \Fil^1 V)$. Since this holds for all $L \in E(K_\infty/K)$, we have $c \in H^1_{\Iw}(K_\infty/K,\Fil^1 V)$.

\section{Proof of the theorem}
\label{secproof}

The only additional complication that arises when proving Theorem \ref{theoB} in full generality is to take care of the possible finite-dimensional quotient in case (2) of Theorem \ref{theoA}. The main result that we need in order to prove Theorem \ref{theoB} is the following (Prop II.3.1 of \cite{Ber05}; recall that $H^1_{\Iw}(K_\cyc/K,V) = \dfont^\dagger(V)^{\psi=1}$ in the notation of ibid).

\begin{prop}
\label{prop1}
If $k \geq 0$ and if $V$ is a de Rham representation of $G_K$ whose Hodge--Tate weights are all $\geq -k$, and which has no subrepresentation all of whose Hodge--Tate weights are $\geq 1-k$, then $H^1_{\Iw,g}(K_{\cyc}/K,V) \subset V^{H_K}$.
\end{prop}

\begin{coro}
\label{coro1}
If $k \geq 0$ and if $V$ is a de Rham representation of $G_K$ whose Hodge--Tate weights are all $\geq -k$, and which has no subrepresentation all of whose Hodge--Tate weights are $\geq 1-k$, then $H^1_{\Iw,g}(K_\infty/K,V)$ is a finite-dimensional $\Qp$-vector space of dimension $\leq \dim(V)$.
\end{coro}

\begin{proof}
We have $H^1_{\Iw,g}(K_\infty/K,V) \subset \projlim_{L \in E(K_\infty/K)} H^1_{\Iw,g}(L_{\cyc}/L,V)$ and each of the $H^1_{\Iw,g}(L_{\cyc}/L,V)$ is a finite-dimensional $\Qp$-vector space of dimension $\leq \dim(V)$ by Proposition \ref{prop1}. 
This implies the result by Lemma \ref{lemm1} below.
\end{proof}

\begin{lemm}
\label{lemm1}
If $d \geq 1$ and $\{ V_i \}_{i \in I }$ is a directed inverse system 
of $\Qp$-vector spaces with $\dim V_i \leq d$ for all $i \in I$, then $\dim \projlim_{i \in I} V_i \leq d$.
\end{lemm}

\begin{proof}
Take $c_1, \hdots, c_{d+1} \in \projlim_{i \in I} V_i$. Let $R_i \subset \Qp^{d+1}$ be the set of $(a_1,\hdots,a_{d+1})$ such that $\sum_{n=1}^{d+1} a_n  \cdot \pr_{V_i}(c_n) = 0$ in $V_i$. Each $R_i$ is nonzero since $\dim V_i \leq d$ and by directedness, each finite intersection of $R_i$ is nonzero. Since $\Qp^{d+1}$ is finite-dimensional, the intersection of the $R_i$  is nonzero. This gives a nontrivial linear relation between the $c_n$. 
\end{proof}

Theorem \ref{theoB} now follows from the same d\'evissage arguments as in \cite{Ber05}, which we now give. Let $\Fil^j V$ denote the largest subrepresentation of $V$ whose Hodge--Tate weights are all $\geq j$. Note that if $L/K$ is a finite extension and $V$ is a representation of $G_K$, then the largest $G_L$-subrepresentation $\Fil_L^j V$ of $V$ whose Hodge--Tate weights are all $\geq j$ is stable under $G_K$. Hence it is equal to $\Fil_K^j V$, so we can drop the subscript from the notation. Now take $k \geq 0$. We have an exact sequence
\[ 0 \to H^1_{\Iw,g}(K_\infty/K,\Fil^{1-k} V) \to H^1_{\Iw,g}(K_\infty/K, \Fil^{-k} V) \to H^1_{\Iw,g}(K_\infty/K,\Fil^{-k} V / \Fil^{1-k} V). \]
By Lemma II.3.3 of \cite{Ber05}, we have $\Fil^j(V/\Fil^j V) = \{0\}$ for all $j \in \ZZ$. 
Therefore the quotient $\Fil^{-k} V / \Fil^{1-k} V$ has no subrepresentation all of whose Hodge--Tate weights are $\geq 1-k$. Corollary \ref{coro1} applies, so that $H^1_{\Iw,g}(K_\infty/K,\Fil^{-k} V / \Fil^{1-k} V)$ is a finite-dimensional $\Qp$-vector space, and therefore so is $H^1_{\Iw,g}(K_\infty/K, \Fil^{-k} V) / H^1_{\Iw,g}(K_\infty/K, \Fil^{1-k} V)$. 
Since $V = \Fil^{-k} V$ for $k \gg 0$, it follows by d\'evissage that 
\[ H^1_{\Iw,g}(K_\infty/K,V) / H^1_{\Iw,g}(K_\infty/K, \Fil^1 V) \] 
is a finite-dimensional $\Qp$-vector space. This proves Theorem \ref{theoB}.

\section{The twisted cyclotomic case}

The analogues of Theorem \ref{theoA} and Proposition \ref{prop1} still hold if one replaces the cyclotomic extension of $K$ by a twisted cyclotomic extension, namely (see \S 8 of \cite{Ber16}) if one replaces $K_{\cyc}$ by $K_{\cyc}^{\eta} = (K^{\alg})^{\ker \chi_{\cyc} \cdot \eta}$ with $\eta : G_K \to \Zp^\times$ an unramified character. The corresponding analogue of Theorem \ref{theoB} then also holds.

\begin{theo}
\label{theoC}
If $V$ is a de Rham representation of $G_K$ and if $K_\infty$ contains a twisted cyclotomic extension $K_{\cyc}^\eta$ of $K$, then $H^1_{\Iw,g}(K_\infty/K,V)  / H^1_{\Iw}(K_\infty/K,\Fil^1 V)$ is a finite-dimensional $\Qp$-vector space.
\end{theo}

\begin{proof}
The proof of this result is the same as the proof of Theorem \ref{theoB}, if we have an analogue of Proposition \ref{prop1} for $K_{\cyc}^{\eta}$. We now explain why such an analogue holds. If $P$ is the completion of the maximal unramified extension of $K$, then $P_{\cyc} = P^\eta_{\cyc}$. Take $V$ as in the statement of Proposition \ref{prop1}. The restriction-to-inertia map gives rise to an injective map $H^1_{\Iw}(K_{\cyc}^\eta/K,V) \to H^1_{\Iw}(P_{\cyc}/P,V)$. The space $H^1_{\Iw,g}(K_{\cyc}^\eta/K,V)$ therefore injects into $H^1_{\Iw,g}(P_{\cyc}/P,V)$, and Proposition \ref{prop1} implies that $H^1_{\Iw,g}(P_{\cyc}/P,V)$ is a finite-dimensional $\Qp$-vector space of dimension $\leq \dim(V)$. This now implies the corresponding analogue of Corollary \ref{coro1}, and the rest of the proof is then the same.
\end{proof}

In particular, one can take for $K_\infty$ any extension containing the torsion points of a Lubin--Tate formal group. Indeed, if that group is attached to a uniformizer $\pi$ of a finite extension $F$ of $\Qp$ contained in $K$, and if $\chi_{\pi} : G_F \to \mathcal{O}_F^\times$ is the corresponding character, then $\Nm_{F/\Qp} (\chi_{\pi})$ is of the form $\chi_{\cyc} \cdot \eta$ as above.

\vspace{\baselineskip}

\noindent\textbf{Acknowledgements.} I thank Rustam Steingart for his input, in particular for requesting the twisted cyclotomic case.

\vspace{\baselineskip}

\noindent\textbf{Tool and computational resource disclosure.} ChatGPT was used to improve the style of writing as well as for assistance with the proof of some technical results.

\bibliographystyle{smfalpha}
\bibliography{UnivNorms}
\end{document}
